\documentclass[../DoAn.tex]{subfiles}
\begin{document}

\section{Kết luận}

Trong khuôn khổ đồ án tốt nghiệp, sinh viên đã xây dựng thành công hệ thống RAG (Retrieval-Augmented Generation) Chatbot hỗ trợ sinh viên hỏi-đáp tự động về các quy định, quy chế đào tạo của Đại học Bách khoa Hà Nội. Đồ án đã giải quyết được các vấn đề chính sau:

\begin{itemize}
    \item Xây dựng pipeline xử lý PDF đa phương thức với Docling, PyMuPDF4LLM và olmOCR, đặc biệt giải quyết vấn đề font encoding cho văn bản tiếng Việt;
    
    \item Phát triển OlmOcrLegalChunker với khả năng nhận diện cấu trúc Chương-Điều-Khoản-Điểm, bảo toàn metadata và ngữ cảnh pháp lý;
    
    \item Triển khai EmbeddingPipeline với mô hình E5 Multilingual (1024 chiều) và FaissVectorStore cho semantic search hiệu quả;
    
    \item Xây dựng GeminiRAG tích hợp Gemini 2.5 Flash để sinh câu trả lời chính xác với trích dẫn nguồn;
    
    \item Phát triển Backend API với FastAPI và Frontend với React 18 + TypeScript + Vite.
\end{itemize}

Hệ thống đã xử lý thành công 16 văn bản quy chế đào tạo, tạo ra 1,247 chunks được lưu trong vector store. Kết quả đạt được:
\begin{itemize}
    \item Độ chính xác xử lý tiếng Việt: 99.5\%+
    \item Thời gian phản hồi: 2-4 giây/query
    \item olmOCR success rate: 100\%
\end{itemize}

Tuy nhiên, đồ án vẫn còn tồn tại một số hạn chế:

\begin{itemize}
    \item Chưa có query caching, mỗi query phải embedding lại;
    \item Single-stage retrieval, chưa có re-ranking với cross-encoder;
    \item Chưa tận dụng lịch sử hội thoại, không hỗ trợ follow-up questions;
    \item FAISS Flat index chưa scale được cho dataset rất lớn;
    \item Chưa có cơ chế update tài liệu tự động.
\end{itemize}

\section{Hướng phát triển trong tương lai}

Dựa trên kết quả đạt được và các hạn chế còn tồn tại, các hướng phát triển khả thi trong tương lai bao gồm:

\subsection{Priority 1 - Critical (1-2 tháng)}
\begin{itemize}
    \item Implement query caching (semantic + exact match) để giảm latency
    \item Add conversation history management để hỗ trợ follow-up questions
    \item Build evaluation dataset với 100+ test queries
    \item Implement basic metrics dashboard để monitor hệ thống
\end{itemize}

\subsection{Priority 2 - Important (2-3 tháng)}
\begin{itemize}
    \item Hybrid search kết hợp BM25 + semantic retrieval
    \item Query rewriting với context để cải thiện retrieval
    \item Re-ranking với cross-encoder để tăng precision
    \item Table-aware chunking strategy cho các bảng phức tạp
\end{itemize}

\subsection{Priority 3 - Nice to have (3-6 tháng)}
\begin{itemize}
    \item Web crawler tự động thu thập documents mới từ HUST
    \item Admin dashboard để monitor logs và metrics
    \item User feedback mechanism để cải thiện chất lượng liên tục
    \item Multi-document reasoning cho câu hỏi phức tạp
\end{itemize}

\subsection{Các cải tiến kỹ thuật}
\begin{itemize}
    \item Upgrade FAISS index từ Flat sang IVF hoặc HNSW cho scale lớn hơn
    \item Thử nghiệm các embedding models khác như BGE, Cohere
    \item Tích hợp streaming response để cải thiện UX
    \item Deploy production với Docker và Kubernetes
\end{itemize}

Tổng kết lại, đồ án là bước khởi đầu tốt cho việc ứng dụng công nghệ RAG trong môi trường giáo dục đại học. Hệ thống đã chứng minh tính khả thi và hiệu quả của việc sử dụng AI để hỗ trợ sinh viên tra cứu quy chế đào tạo, mở ra nhiều tiềm năng ứng dụng thực tiễn tại ĐHBK Hà Nội và các trường đại học khác.

\end{document}
